\section{第12章：气体动理论（12.5练习题）}
\sourceproblem{选择1}{218}{224}
\problemtext 给定理想气体p、V、T，求分子数。
\workings $pV=NkT$，所以 $N=pV/(kT)$，选B。
\crosscheck $pV/(RT)$ 为物质的量而非分子数。
\sourceproblem{选择2}{219}{225}
\problemtext 标准状态下每立方米理想气体分子数。
\workings $n=p/(kT)=1.013\times10^5/(1.38\times10^{-23}\cdot273)\simeq2.69\times10^{25}$ 每立方米，选C。
\crosscheck 采用原书标准状态1 atm、273 K。
\sourceproblem{选择3}{219}{225}
\problemtext 273 K、$1.013\times10^5$ Pa、密度1.25 kg每立方米，判断气体。
\workings $M=\rho RT/p\simeq0.0280$ kg/mol，氮气或一氧化碳，选C。
\crosscheck 摩尔质量不能唯一鉴别等摩尔质量的气体。
\sourceproblem{选择4}{219}{225}
\problemtext 平衡理想气体速度x分量平方的平均值。
\workings 各向同性给 $\overline{v_x^2}=\overline{v^2}/3=kT/m$，选D。
\crosscheck 平均速度分量为零，不代表平方平均为零。
\sourceproblem{选择5}{219}{225}
\problemtext 同种气体数密度相同，方均根速率比1:2:4，求压强比。
\workings $p=nm\overline{v^2}/3$，故1:4:16，选C。
\crosscheck 压强正比于方均根速率的平方。
\sourceproblem{选择6}{219}{225}
\problemtext 同温氢气和氧气的压强、密度及分子速率判断。
\workings $v_{\rm rms}=\sqrt{3kT/m}$，氢气方均根速率较大，选D。
\crosscheck 个别分子速率无固定大小顺序；压强还需数密度。
\sourceproblem{选择7}{219}{225}
\problemtext 判断关于温度统计意义的四项陈述。
\workings $\overline E_{\rm tr}=3kT/2$，温度为统计量，反映热运动程度；前三项成立，第四项把温度赋予单个分子不当，选B。
\crosscheck 限于题设平衡热运动语境。
\sourceproblem{选择8}{219}{225}
\problemtext 氦、氮气质量密度和平均平动动能相同，比较温度和压强。
\workings 同温且 $p=\rho RT/M$，氦摩尔质量较小，压强较大，选C。
\crosscheck 质量密度相同不代表数密度相同。
\sourceproblem{选择9}{219}{225}
\problemtext 同温同压氦和氧气平均能量及平均平动动能比较（选项续226页）。
\workings 平动均为 $3kT/2$；经典刚性分子模型下总平均热运动能分别为 $3kT/2,5kT/2$，选C。
\crosscheck 忽略振动激发及内能参考常数。
\sourceproblem{选择10}{220}{226}
\problemtext 1 mol刚性双原子理想气体内能。
\workings 五个二次自由度给 $U=5RT/2$，选C。
\crosscheck R用于摩尔量，k用于单分子。
\sourceproblem{选择11}{220}{226}
\problemtext 标准状态氧气与氦气体积比1:2，求内能比。
\workings 同温同压物质的量比1:2，故 $U_1/U_2=(5/3)(1/2)=5/6$，选C。
\crosscheck 氧气按刚性双原子处理。
\sourceproblem{选择12}{220}{226}
\problemtext p、V已知的刚性氢气内能。
\workings $U=5\nu RT/2=5pV/2$，选A。
\crosscheck 经典刚性双原子模型。
\sourceproblem{选择13}{220}{226}
\problemtext 同温氢氧速率曲线a、b及最概然速率比。
\workings $v_p=\sqrt{2kT/m}$，氧与氢比 $\sqrt{2/32}=1/4$。原图a峰位较左，故a为氧气，选B。
\crosscheck 重分子分布较窄且峰较高，面积仍为1。
\sourceproblem{选择14}{220}{226}
\problemtext 比较最概然、平均、方均根速率。
\workings 三者分别为 $\sqrt{2kT/m},\sqrt{8kT/(\pi m)},\sqrt{3kT/m}$，严格递增，选C。
\crosscheck T取正值。
\sourceproblem{选择15}{220}{226}
\problemtext 同种气体T1大于T2，比较最概然速率及分布峰高。
\workings $v_p\propto\sqrt T$，麦克斯韦分布 $f(v_p)=4/(\sqrt\pi e v_p)$，峰高与温度平方根成反比，选B。
\crosscheck 此处峰高随温度升高而减小，并非分子总数降低。
\sourceproblem{选择16}{220}{226}
\problemtext 四幅曲线中选择同温氮气与氦气分布。
\workings 氮气分子较重，峰更靠左且高；氦峰靠右且低，均归一化。原图B满足定性关系，选B。
\crosscheck 示意图不用于精确量取峰位比。
\sourceproblem{选择17}{221}{227}
\problemtext 封闭定容气体平均速率增为2倍，比较温度和压强。
\workings $\bar v\propto\sqrt T$，故T为4倍；$p=nkT$，p也为4倍，选D。
\crosscheck 分子数与容积固定。
\sourceproblem{选择18}{221}{227}
\problemtext 不同气体最概然速率相等，比较其平均和方均根速率。
\workings 三种特征速率均只依赖 $T/m$，所以均相等，选A。
\crosscheck 平衡麦克斯韦模型。
\sourceproblem{选择19}{221}{227}
\problemtext 氧气绝对温度倍增且完全解离，求氧原子平均速率比。
\workings 质量减半、温度倍增，$\bar v'/\bar v=\sqrt{2/(1/2)}=2$，选B。
\crosscheck 给定末态温度，解离能不另用于求温度。
\sourceproblem{选择20}{221}{227}
\problemtext 麦克斯韦曲线被v0分成等面积两部分，解释v0。
\workings $\int_0^{v_0} f(v)\,dv=\int_{v_0}^{\infty}f(v)\,dv=1/2$，v0是速率中位数，选D。
\crosscheck 不等于最概然速率或平均速率的一般定义。
\sourceproblem{选择21}{221}{227}
\problemtext 速率分布函数的物理意义。
\workings $f(v)dv=dN/N$，表示每单位速率间隔的概率密度，选B。
\crosscheck 严格为密度而非单点概率；百分比需乘100%。
\sourceproblem{选择22}{221}{227}
\problemtext 解释指定速率区间内平动动能积分。
\workings $Nf(v)dv$ 为该微元区间分子数，$\int_{v_1}^{v_2}mv^2Nf(v)dv/2$ 为区间分子平动动能总和，选D。
\crosscheck 求区间平均须再除该区间分子数。
\sourceproblem{选择23}{221}{227}
\problemtext 求速率位于给定区间的分子平均速率。
\workings 条件平均 $\bar v_{\rm int}=\int_{v_1}^{v_2}vf(v)dv/\int_{v_1}^{v_2}f(v)dv$，选C。
\crosscheck 要求分母非零。
\sourceproblem{选择24}{221}{227}
\problemtext 总分子数N及分布函数已知，求区间平均速率（选项续228页）。
\workings 分子数权重N在分子分母抵消，$\bar v_{\rm int}=\int_{v_1}^{v_2}vf(v)dv/\int_{v_1}^{v_2}f(v)dv$，选B。
\crosscheck 与选择23同一概念，保留原题号。
\sourceproblem{选择25}{222}{228}
\problemtext 定温膨胀时平均碰撞频率及平均自由程变化。
\workings 硬球模型 $\ell=1/(\sqrt2\pi d^2n)$，$Z=\bar v/\ell$。n下降，故自由程增大、频率减小，选B。
\crosscheck 有效碰撞直径恒定。
\sourceproblem{选择26}{222}{228}
\problemtext 封闭容器1 mol氦气，平均自由程由哪个状态量决定。
\workings $\ell=V/(\sqrt2\pi d^2N_A)$，固定碰撞直径时仅取决于V，选B。
\crosscheck 真实有效截面可随温度变化，教材为硬球模型。
\sourceproblem{选择27}{222}{228}
\problemtext 定容升温时自由程与碰撞频率变化。
\workings n固定则自由程不变，$Z\propto\sqrt T$，选D。
\crosscheck 有效碰撞直径不变。
\sourceproblem{选择28}{222}{228}
\problemtext 定容温度减半且碰撞球半径不变，求自由程。
\workings $\ell=1/(\sqrt2\pi d^2n)$ 不变，选B。
\crosscheck 温度不直接进入固定数密度的硬球自由程。
\sourceproblem{选择29}{222}{228}
\problemtext 定容绝对温度变为1/4，求自由程及碰撞频率。
\workings $\ell=\ell_0,\ Z=Z_0/2$，选B。
\crosscheck 保持粒子数及碰撞直径。
\sourceproblem{选择30}{222}{228}
\problemtext 恒压下平均碰撞频率随温度的关系。
\workings $n=p/(kT)$，$\ell\propto T$，$\bar v\propto\sqrt T$，故 $Z\propto T^{-1/2}$，选C。
\crosscheck 固定有效碰撞直径及气体种类。
\sourceproblem{填空1}{222}{228}
\problemtext 电子管压强1.0乘10的负5次方mmHg，27摄氏度，求数密度。
\workings $p=10^{-5}(1.013\times10^5/760)$ Pa，$n=p/(kT)\simeq3.22\times10^{17}$ 每立方米。
\crosscheck 采用300 K，毫米汞柱换算不可漏掉760。
\sourceproblem{填空2}{222}{228}
\problemtext 等温压缩，初始数密度1.96乘10的24次方，压强倍增。
\workings $n=p/(kT)$，故 $n'=3.92\times10^{24}$ 每立方米。
\crosscheck 物质的量固定。
\sourceproblem{填空3}{222}{228}
\problemtext 判断三种理想气体微分关系（第三式续229页）。
\workings 固定物质的量 $\nu$ 时 $pdV+Vdp=\nu RdT$。若 $pdV=\nu RdT$，则等压；若 $Vdp=\nu RdT$，则等容；若 $pdV+Vdp=0$，则等温。
\crosscheck 假设关系在整个准静态过程成立。
\sourceproblem{填空4}{223}{229}
\problemtext 分子物理学的研究对象及基本方法。
\workings 从物质微观分子运动和相互作用出发研究宏观热现象及热性质，基本方法为统计方法。
\crosscheck 宏观量与大量微观粒子的统计平均相联系。
\sourceproblem{填空5}{223}{229}
\problemtext 理想气体分子模型主要三点。
\workings 分子自身大小相对间距可忽略；除碰撞外分子间相互作用可忽略；分子之间及与器壁的碰撞视为完全弹性，碰撞时间相对自由运动时间可忽略。
\crosscheck 各向同性等是平衡态统计假设，不是单个分子的运动限制。
\sourceproblem{填空6}{223}{229}
\problemtext 方均根速率450 m/s、压强7万Pa，求质量密度。
\workings $\rho=3p/v_{\rm rms}^2=210000/202500\simeq1.037$ kg每立方米。
\crosscheck 使用方均根速率而非平均速率。
\sourceproblem{填空7}{223}{229}
\problemtext 固定质量，等压升温及等温增压时密度变化。
\workings $\rho=pM/(RT)$，等压升温密度减小，等温增压密度增大。
\crosscheck 固定气体组成。
\sourceproblem{填空8}{223}{229}
\problemtext 容积0.01立方米，气体质量100 g，方均根速率200 m/s，求压强。
\workings $\rho=0.1/0.01=10$ kg每立方米，$p=\rho v_{\rm rms}^2/3\simeq1.333\times10^5$ Pa。
\crosscheck 克换千克。
\sourceproblem{填空9}{223}{229}
\problemtext 273 K、压强0.01 atm、密度0.0124 kg每立方米，求方均根速率。
\workings $v_{\rm rms}=\sqrt{3p/\rho}=\sqrt{3(1013)/0.0124}\simeq495$ m/s。
\crosscheck 温度不是此组p、rho数据求速率所需的独立量。
\sourceproblem{填空10}{223}{229}
\problemtext 微观压强公式中的两个决定因素。
\workings $p=2n\overline E_{\rm tr}/3$，决定于分子数密度及平均平动动能。
\crosscheck 对各向同性平衡理想气体。
\sourceproblem{填空11}{223}{229}
\problemtext 三个气体数密度比4:2:1，平均平动动能比1:2:4，求压强比。
\workings $p=2n\overline E_{\rm tr}/3$，故 $p_A:p_B:p_C=1:1:1$。
\crosscheck 原书末尾重复pA，按三容器顺序理解。
\sourceproblem{填空12}{223}{229}
\problemtext 1 mol刚性氧气27摄氏度，求内能、平均平动动能、平均能量、平均动能。
\workings $U=5RT/2=6232.5$ J；$\overline E_{\rm tr}=3kT/2=6.21\times10^{-21}$ J；经典刚性模型下平均能量及平均动能均 $5kT/2=1.035\times10^{-20}$ J。
\crosscheck 内能取热运动能零点，忽略振动及相互作用势能。
\sourceproblem{填空13}{223}{229}
\problemtext 质量m'的刚性双原子氢气，求平均平动动能、平均动能及内能。
\workings 依次为 $3kT/2,5kT/2,5m'RT/(2M)$。
\crosscheck M为氢气摩尔质量。
\sourceproblem{填空14}{223}{229}
\problemtext 能量10的12次方eV完全被0.1 mol氖气吸收，求升温。
\workings 若容器定容且无其他能量损失，$Q=1.6\times10^{-7}$ J，$\Delta T=Q/(3\nu R/2)\simeq1.28\times10^{-7}$ K。
\crosscheck 若气体对外做功，不能仅凭吸收能量确定升温；按氖管定容理解。
\sourceproblem{填空15}{223}{229}
\problemtext 氮气压强1.33 Pa、7摄氏度，求数密度及质量密度。
\workings $n=p/(kT)\simeq3.44\times10^{20}$ 每立方米；$\rho=pM/(RT)\simeq1.60\times10^{-5}$ kg每立方米。
\crosscheck T取280 K，1 mol信息不影响密度。
\sourceproblem{填空16}{223}{229}
\problemtext 2 g氢和2 g氦同温同容，求平动动能、压强、内能比。
\workings 分子平均平动能比1:1；物质的量为1 mol和0.5 mol，压强比2:1；内能比 $(5/2)/(3/4)=10:3$。
\crosscheck 氢视为刚性双原子，氦为单原子。
\sourceproblem{填空17}{224}{230}
\problemtext 平均平动动能与温度及其统计意义。
\workings $\overline E_{\rm tr}=3kT/2$，温度是大量分子平均平动动能的量度。
\crosscheck 单个分子瞬时动能不是热力学温度。
\sourceproblem{填空18}{224}{230}
\problemtext 平均平动动能1.06乘10的负19次方J，求温度。
\workings $T=2\overline E_{\rm tr}/(3k)\simeq5.12\times10^3$ K。
\crosscheck 使用题给k；高温可能有其他自由度，但平动关系仍成立于经典平衡模型。
\sourceproblem{填空19}{224}{230}
\problemtext 单原子理想气体E=3RT/2的物理意义。
\workings 为1 mol单原子理想气体的热运动内能，亦即摩尔内能；一般 $U=3\nu RT/2$。
\crosscheck 不是单分子平均能量。
\sourceproblem{填空20}{224}{230}
\problemtext 1 mol单原子气体恒压由0摄氏度加热到100摄氏度，求内能增量。
\workings $\Delta U=3\nu R\Delta T/2=1246.5$ J。
\crosscheck 不是恒压吸热量 $5
u RDelta T/2$。
\sourceproblem{填空21}{224}{230}
\problemtext 同温氦、氖、氩三曲线，辨别a与c。
\workings 原图a峰最靠左，c最靠右；$v_p\propto m^{-1/2}$，故a为氩，c为氦，b为氖。
\crosscheck 峰高随质量平方根增加。
\sourceproblem{填空22}{224}{230}
\problemtext 说明最概然速率以上概率积分及全速率平动动能积分。
\workings $\int_{v_p}^\infty f(v)dv$ 为速率大于最概然速率的分子比例；$\int_0^\infty mv^2 f(v)dv/2$ 为单分子平均平动动能。
\crosscheck 两积分量纲分别为1及J。
\sourceproblem{填空23}{224}{230}
\problemtext 表达速率大于v0的分子数、条件平均速率及概率。
\workings 依次为 $N\int_{v_0}^\infty f(v)dv$、$\int_{v_0}^\infty vf(v)dv/\int_{v_0}^\infty f(v)dv$、$\int_{v_0}^\infty f(v)dv$。
\crosscheck 第三问的时间抽样概率按平衡平稳且遍历的统计模型理解。
\sourceproblem{填空24}{224}{230}
\problemtext 同温氢、氦速率分布图，左峰标1000 m/s，求各最概然速率。
\workings 氦质量较大，左峰为氦，$v_{p,\rm He}=1000$ m/s；氢为右峰，$v_{p,\rm H_2}=\sqrt2\,1000$ m/s。
\crosscheck 图中右峰无数字，应由分子质量比推算。
\sourceproblem{填空25}{224}{230}
\problemtext 氢、氧平均平动动能相同，方均根速率是否相同。
\workings 不同；$v_{\rm rms}=\sqrt{2\overline E_{\rm tr}/m}$，氢为氧的4倍。
\crosscheck 使用分子而非原子质量。
\sourceproblem{填空26}{224}{230}
\problemtext 给定分布函数，求速率在最概然速率至无穷范围的概率。
\workings $\Delta N/N=\int_{v_p}^\infty f(v)dv$。若代入麦克斯韦分布，结果为 $\operatorname{erfc}(1)+2/(e\sqrt\pi)\simeq0.5724$。
\crosscheck 最概然速率不是中位数。
\sourceproblem{填空27}{224}{230}
\problemtext 同温氢与氧平均平动动能比及方均根速率比。
\workings 能量比1:1，速率比 $\sqrt{32/2}:1=4:1$。
\crosscheck 同温不等于同速率。
\sourceproblem{填空28}{224}{230}
\problemtext 等压由V0膨胀至2V0，求自由程、平均速率、平动动能比（续231页）。
\workings $T/T_0=2,\ n/n_0=1/2$，故三比值依次为 $2,\sqrt2,2$。
\crosscheck 硬球直径不变，分子数固定。
\sourceproblem{填空29}{225}{231}
\problemtext 先定容升温倍增，再等温体积倍增，求自由程倍数。
\workings 固定N及碰撞直径，$\ell\propto V$，故最终为2倍。
\crosscheck 第一阶段不改变自由程。
\sourceproblem{填空30}{225}{231}
\problemtext 先定容升温倍增，再等压体积倍增，求自由程倍数。
\workings 末态体积为初态2倍，故 $\ell/\ell_0=2$。
\crosscheck 末态温度为初态4倍，但固定碰撞截面的自由程仅由数密度决定。

\sourceproblem{计算1}{225}{231}
\problemtext 5 L容器内含 $10^{15}$ 个氧分子、$4\times10^{15}$ 个氮分子及 $3.3\times10^{-7}$ g氩，333 K，求混合压强。
\workings 取 $M_{\rm Ar}=40$ g/mol、$N_A=6.02\times10^{23}$ 每摩尔，氩原子数 $N_{\rm Ar}=3.3\times10^{-7}N_A/40=4.9665\times10^{15}$。由各组分分压之和，
\[
p=\frac{(N_{\rm O_2}+N_{\rm N_2}+N_{\rm Ar})kT}{V}
\simeq9.16\times10^{-3}\ \mathrm{Pa}.
\]
\crosscheck 体积为 $5\times10^{-3}$ 立方米；氩为单原子气体，压强与内自由度无关。
\sourceproblem{计算2}{225}{231}
\problemtext 氧气压强 $1.01\times10^5$ Pa、温度27摄氏度，求数密度、质量密度、平均平动动能和平均能量。
\workings 取T=300 K、氧气摩尔质量0.032 kg/mol，有
\[
n=\frac p{kT}\simeq2.44\times10^{25}\ \mathrm{m}^{-3},\quad
\rho=\frac{pM}{RT}\simeq1.296\ \mathrm{kg/m^3}.
\]
$\overline E_{\rm tr}=3kT/2=6.21\times10^{-21}$ J；按常温刚性双原子模型，$\overline E=5kT/2=1.035\times10^{-20}$ J。
\crosscheck 平均总热运动能含转动项；原题未指定振动自由度，采用本章常温刚性模型并取热运动能零点。
\sourceproblem{计算3}{225}{231}
\problemtext 2.56 kg氧气的总平动动能为 $4\times10^5$ J，求分子平均平动动能及温度。
\workings 物质的量为80 mol，$N=80(6.02\times10^{23})$。于是
\[
\overline E_{\rm tr}=\frac{4\times10^5}{N}\simeq8.31\times10^{-21}\ \mathrm J,\qquad
T=\frac{2\overline E_{\rm tr}}{3k}\simeq401\ \mathrm K.
\]
\crosscheck 亦可用 $E_{\rm tr}=3\nu RT/2$；不能以双原子总内能的5个自由度替换题给平动能。
\sourceproblem{计算4}{225}{231}
\problemtext 2立方米容器内有 $10^{25}$ 个氢分子和 $4\times10^{25}$ 个氧分子，400 K，求分子总动能及压强。
\workings 两气体按刚性双原子模型，$N=5\times10^{25}$，总热运动动能 $E=5NkT/2=6.90\times10^5$ J，压强 $p=NkT/V=1.38\times10^5$ Pa。若仅指质心平动动能，则为 $3NkT/2=4.14\times10^5$ J。
\crosscheck 原题使用动能总和而非平动动能，主解计转动；压强由平动动能决定，$pV=2E_{\rm tr}/3$。
\sourceproblem{计算5}{226}{232}
\problemtext 0.1 kg氧气在10 atm、47摄氏度的刚性容器内，漏气后压强为原来的5/8，温度27摄氏度，求容积和漏出质量。
\workings 取 $p_1=1.013\times10^6$ Pa，$T_1=320$ K、$T_2=300$ K。初态
\[
V=\frac{m_1 RT_1}{Mp_1}\simeq8.20\times10^{-3}\ \mathrm{m^3}.
\]
定容但物质的量变化，$m_2/m_1=(p_2/p_1)(T_1/T_2)=2/3$，故漏出质量 $m_1-m_2=0.03333$ kg。
\crosscheck 压强比不能直接当作剩余质量比，需计降温；无需假设漏气过程可逆。
\sourceproblem{计算6}{226}{232}
\problemtext 刚性双原子气体体积 $2\times10^{-3}$ 立方米、内能675 J、分子数 $5.4\times10^{22}$，求压强、温度及平均平动动能。
\workings $U=5NkT/2=5pV/2$，故 $p=2U/(5V)=1.35\times10^5$ Pa。$T=2U/(5Nk)\simeq362.3$ K；$\overline E_{\rm tr}=3U/(5N)=7.50\times10^{-21}$ J。
\crosscheck 平动能总和为405 J，占刚性模型内能的3/5。
\sourceproblem{计算7}{226}{232}
\problemtext 氢、氧同温，氢分子平均平动动能 $6.21\times10^{-21}$ J，求氧气温度、平均平动动能和方均根速率。
\workings 由 $T=2\overline E_{\rm tr}/(3k)=300$ K，氧分子平均平动动能相同。取 $m_{\rm O_2}=0.032/N_A$ kg，
\[
v_{\rm rms}=\sqrt{2\overline E_{\rm tr}/m_{\rm O_2}}\simeq483.5\ \mathrm{m/s}.
\]
\crosscheck 同温同平均平动能不意味着同速率。
\sourceproblem{计算8}{226}{232}
\problemtext 1 mol氧气装于1立方米容器，以10 m/s整体运动，突然停止后定向动能的80%转为热运动动能，求升温、增压。
\workings 假定氧气重新达到平衡、体积固定，按刚性双原子模型，热运动内能增加 $\Delta U=0.8Mv^2/2=1.28$ J。于是 $\Delta T=\Delta U/(5R/2)\simeq0.0616$ K，$\Delta p=R\Delta T/V=0.512$ Pa。
\crosscheck 这里的热运动动能按平动和转动之和理解；若仅问未发生转动均分的瞬时平动态，则用3个自由度会给不同值，不能当作最终平衡温度。
\sourceproblem{计算9}{226}{232}
\problemtext 氦原子方均根速率550 m/s时求温度；另求5800 K时的方均根速率。
\workings 取氦摩尔质量0.004 kg/mol，$T=Mv_{\rm rms}^2/(3R)\simeq48.54$ K；$v_{\rm rms}'=\sqrt{3R(5800)/M}\simeq6.01\times10^3$ m/s。
\crosscheck 后者按经典理想氦原子模型；实际太阳气体可部分电离，不能据此断言全部为中性氦。
\sourceproblem{计算10}{227}{233}
\problemtext 373 K、压强 $1.01\times10^3$ Pa、密度 $1.24\times10^{-5}$ 克每立方厘米，求方均根速率、摩尔质量及可能气体。
\workings 密度换为0.0124 kg每立方米，故
\[
v_{\rm rms}=\sqrt{3p/\rho}\simeq494.3\ \mathrm{m/s},\quad
M=\rho RT/p\simeq0.0381\ \mathrm{kg/mol}.
\]
与常见氩气摩尔质量约0.040 kg/mol较接近，教材常见气体候选中可推测为氩。
\crosscheck 题给数据对应约38.1 g/mol，与40 g/mol存在约5\%差异；不将该推测写作唯一鉴定，混合气体或其他组成也可能有同一平均摩尔质量。
\sourceproblem{计算11}{227}{233}
\problemtext 平衡混合物中两种单原子气体质量分别 $m_1,m_2$，各自内能均为E，求平均速率比和整体平均速率。
\workings 同温且 $E=3N_i kT/2$，所以两组分粒子数相同。对每组分 $\bar v_i=\sqrt{8kT/(\pi m_i/N_i)}=\sqrt{16E/(3\pi m_i)}$。于是
\[
\frac{\bar v_1}{\bar v_2}=\sqrt{\frac{m_2}{m_1}},\qquad
\bar v=\frac{\bar v_1+\bar v_2}{2}
=\sqrt{\frac{4E}{3\pi}}\left(\frac1{\sqrt{m_1}}+\frac1{\sqrt{m_2}}\right).
\]
\crosscheck 混合物的平均速率按粒子数加权而非质量加权；不能用整体均方根替代平均速率。
\sourceproblem{计算12}{227}{233}
\problemtext 固定气体从300 K升至450 K、体积减半，求压强、平均平动动能及方均根速率的变化。
\workings $p_2/p_1=(450/300)/(1/2)=3$；$\overline E_{{\rm tr},2}/\overline E_{{\rm tr},1}=3/2$，增量 $3k(150)/2=3.105\times10^{-21}$ J；$v_{{\rm rms},2}/v_{{\rm rms},1}=\sqrt{3/2}\simeq1.2247$。
\crosscheck 压强增加200\%、平均平动能增加50\%、方均根速率增加约22.47\%；未给分子质量及初压，不能给后二者中速率及压强的绝对增量。
\sourceproblem{计算13}{227}{233}
\problemtext 五组粒子数2、4、6、8、2，速率分别10、20、30、40、50 m/s，求平均及方均根速率。
\workings 总数22，$\sum N_iv_i=700$ m/s，$\sum N_iv_i^2=25000$ 平方米每平方秒，故
\[
\bar v=700/22\simeq31.82\ \mathrm{m/s},\qquad
v_{\rm rms}=\sqrt{25000/22}\simeq33.71\ \mathrm{m/s}.
\]
\crosscheck 采用粒子数权重；方均根不小于平均值。
\sourceproblem{计算14}{227}{233}
\problemtext 速率分布 $f(v)=Av$ 于 $0\le v\le v_c=100$ m/s，其他速率处为零，求A和方均根速率。
\workings 归一化给 $\int_0^{v_c}Av\,dv=Av_c^2/2=1$，故 $A=2/v_c^2=2\times10^{-4}$ 平方秒每平方米。于是
\[
\overline{v^2}=\int_0^{v_c}Av^3dv=v_c^2/2,\qquad v_{\rm rms}=100/\sqrt2\simeq70.71\ \mathrm{m/s}.
\]
\crosscheck 速率定义域非负；N不进入归一化分布的均方根，亦不要求该分布是麦克斯韦分布。
\sourceproblem{计算15}{228}{234}
\problemtext 氢分子有效直径 $2\times10^{-10}$ m，标准状况下求单分子一秒内平均碰撞次数。
\workings 取 $T=273$ K、$p=1.013\times10^5$ Pa，$n=p/(kT)$，$\ell=1/(\sqrt2\pi d^2n)\simeq2.09\times10^{-7}$ m。$\bar v=\sqrt{8RT/(\pi M)}\simeq1.70\times10^3$ m/s，故 $Z=\bar v/\ell\simeq8.12\times10^9$ 每秒，一秒平均碰撞约 $8.1\times10^9$ 次。
\crosscheck 使用平均速率而非方均根速率；硬球碰撞截面为 $\pi d^2$，相对运动带来 $\sqrt2$ 因子。
\sourceproblem{应用1}{228}{234}
\problemtext 将温度 $10^8$ K星球中的质子视为理想气体，求方均根速率和平均平动动能（eV）。
\workings 取质子摩尔质量约 $10^{-3}$ kg/mol，$v_{\rm rms}=\sqrt{3RT/M}\simeq1.58\times10^6$ m/s；$\overline E_{\rm tr}=3kT/2=2.07\times10^{-15}$ J，约 $1.29\times10^4$ eV。
\crosscheck 质子速率约为光速的0.0053，非相对论近似合理；实际等离子体含电子及相互作用，题设只求质子热运动统计量。
\sourceproblem{应用2}{228}{234}
\problemtext 20立方厘米电子管，300 K，压强 $5\times10^{-6}$ mmHg，空气按刚性双原子模型，求分子数及平动、转动、总动能。
\workings $p=5\times10^{-6}(1.013\times10^5/760)\simeq6.6645\times10^{-4}$ Pa，$V=2.0\times10^{-5}$ 立方米，$N=pV/(kT)\simeq3.22\times10^{12}$。
\[
E_{\rm tr}=\tfrac32pV\simeq2.00\times10^{-8}\ \mathrm J,\quad
E_{\rm rot}=pV\simeq1.33\times10^{-8}\ \mathrm J,\quad
E_{\rm tot}=\tfrac52pV\simeq3.33\times10^{-8}\ \mathrm J.
\]
\crosscheck 三者是全体分子的平均能量之和，非单分子能量；稀薄气体按题设温度及平衡分布计算。
\sourceproblem{应用3}{228}{234}
\problemtext 定容自行车内胎7摄氏度时压强 $4.0\times10^5$ Pa，升至37摄氏度，求末压。
\workings 若给定为绝对压强，$p_2=p_1T_2/T_1=4.0\times10^5(310/280)\simeq4.43\times10^5$ Pa。若给定为表压且环境压强 $p_a$ 恒定，则 $p_{g,2}=(p_{g,1}+p_a)310/280-p_a$；例如 $p_a=1.013\times10^5$ Pa时约 $4.54\times10^5$ Pa。
\crosscheck 状态方程必须用绝对压强；题干称空气压强，主解按绝对值，另注明实际胎压表约定。
\sourceproblem{应用4}{229}{235}
\problemtext 装刚性氧气的钢瓶随车以40 m/s运动，突然停止，气体原定向动能的80%转为热运动动能，求升温。
\workings 对质量m气体，$\Delta U=0.8mv^2/2=5mR\Delta T/(2M)$。消去质量，
\[
\Delta T=\frac{0.8Mv^2}{5R}\simeq0.986\ \mathrm K .
\]
\crosscheck 与计算8同一最终平衡模型；此处钢瓶自身的机械动能不计入气体获得的能量。
\sourceproblem{应用5}{229}{235}
\problemtext 50 m深湖水中4摄氏度、体积 $10^{-5}$ 立方米空气泡升至17摄氏度湖面，求体积，给大气压 $1.013\times10^5$ Pa。
\workings 取水密度 $10^3$ kg每立方米、$g=9.8$ m/s平方，忽略表面张力和气体交换。初压 $p_1=p_0+\rho gh=5.913\times10^5$ Pa，末压 $p_2=p_0$。由固定物质的量状态方程，
\[
V_2=V_1\frac{p_1}{p_0}\frac{290}{277}\simeq6.11\times10^{-5}\ \mathrm{m^3}.
\]
\crosscheck 温度须用K；若计曲率压差、气体溶解及水密度变化则需另给参数。
